% Drawing for measurement `jfet-switch`
% Stands for the deck lines:
%   VS s 0 0
%   JW l g s DUT
%   RL l 0 1000
%   RG in g 1000 / VIN in 0 PWL(...)
%   .tran 1e-10 4.2e-6
\documentclass[border=8pt]{standalone}
\usepackage[siunitx,RPvoltages,european]{circuitikz}
\ctikzset{bipoles/length=0.9cm, label/align=straight, inductors/scale=0.8}
\ctikzset{inductor=american}
\usepackage{amsmath}
\begin{document}
\begin{circuitikz}[font=\footnotesize, line width=0.5pt]
\draw (0,0) node[njfet](Q){};
\node[anchor=west] at ($(Q.D)!0.5!(Q.S) + (0.55,0)$) {DUT};
\draw (Q.S) to[short] ++(0,-0.3) to[rmeter, t=A, l={$i_s$}] ++(0,-1.4) to[vsource, l={$v_s$ 0 or \SI{1}{V}}, invert] ++(0,-1.4) node[ground]{};
\draw (Q.D) to[short] ++(0,0.4) to[short] ++(2.4,0) coordinate(l) to[R, l={$R_L$ \SI{1}{k\ohm} / \SI{10}{k\ohm}}] ++(0,-3.2) node[ground]{};
\draw (l) to[short, -o] ++(1,0) node[right]{$v_{load}$};
\draw (Q.G) to[short] ++(-0.4,0) to[R, l={$R_G$ \SI{1}{k\ohm}}] ++(-1.6,0) to[short] ++(-0.8,0) to[vsourcesquare, l_={\shortstack[r]{$0\to-\SI{10}{V}$\\ (or $2V_{GS(\mathrm{off})}$)\\ \SI{10}{ns} edges}}, invert] ++(0,-1.9) node[ground]{};
\end{circuitikz}
\end{document}
